\label{chap:83-09}

% Unidad U012. Biografía estructural completa de pi.

% La biografía de pi que seguía en esta residencia es exactamente la ya
% publicada en el capítulo 8, salvo el namespace de sus etiquetas. Se conserva
% en la fuente por su procedencia, pero no se imprime dos veces.
\iffalse

\section{Construction du caractère de clôture et des trajectoires généalogiques structurelles de \(\pi\)}
\label{p12-83-c9-s1:chap:biografia-pi-autonoma}

\subsection{Séparation causale entre germe, caractère et évaluation}

La construction se divise en trois strates :
\[
 \text{sélection finie}
 \longrightarrow
 \text{construction du caractère de clôture}
 \longrightarrow
 \text{évaluation archimédienne}.
\]
Le développement décimal n’intervient pas dans la première strate. La chaîne des types
est
\begin{equation}
 104\,976
 \longrightarrow468\,\mathcal U_6
 \xrightarrow{\Psi}243\,\mathcal W_6
 \lhook\joinrel\longrightarrow\mathbb F_3^6,
 \qquad|\mathbb F_3^6|=729.
 \label{p12-83-c9-s1:eq:cadena-tipica-pi-autonoma}
\end{equation}
Les cardinalités comptent, respectivement, les conditions initiales
enrichies, les émissions distinctes, les mots ternaires visibles et les éléments
de la fibre ambiante complète.

Pour \(U_6=(u_1,\ldots,u_6)\), la projection est
\[
 \Psi(U_6)=((-u_1)\bmod3,\ldots,(-u_6)\bmod3).
\]
Le recensement des \(43\) orbites de \(D_3\) démontre que l’orbite de
clôture est la seule dont les six orientations possèdent à la fois :
\[
 \text{cinq relèvements},\qquad
 \text{masse }1008,\qquad
 1008=432+4\cdot144.
\]
Le calibre
\[
 \operatorname{diag}_c=6,
 \qquad
 \operatorname{card}=ES
\]
sélectionne un unique représentant :
\begin{equation}
 \boxed{w_6^\pi=010211.}
 \label{p12-83-c9-s1:eq:semilla-pi-autonoma}
\end{equation}
La sélection utilise le catalogue, les multiplicités, l’action orbitale et l’orientation ;
elle ne consulte pas un préfixe de \(\pi\).

\subsection{Relèvements \(L_0,L_1\) et 1ʳᵉ frontière coinductive}

Sur les vecteurs-lignes de \(\mathbb F_3^6\), soient
\[
L_0=
\begin{pmatrix}
2&2&2&1&2&1\\
2&1&2&2&1&1\\
1&0&0&1&0&2\\
0&0&1&1&1&0\\
2&2&2&0&2&1\\
2&1&2&1&0&0
\end{pmatrix},
\quad
L_1=
\begin{pmatrix}
2&1&2&0&2&2\\
1&2&1&1&2&0\\
1&2&1&1&1&2\\
0&1&0&0&2&1\\
2&0&2&1&0&1\\
0&2&2&2&1&1
\end{pmatrix}.
\]
Si \(b_1=w_6^\pi\), la récurrence des blocs est
\[
 b_{k+1}=b_kL_{\varepsilon_k},
 \qquad
 (\varepsilon_1,\varepsilon_2,\varepsilon_3,\varepsilon_4)=(0,0,1,1),
\]
et le mot accumulé est formé par concaténation :
\[
 w_{6(k+1)}=w_{6k}\mathbin{\|}b_{k+1}.
\]
On ne multiplie pas un mot de longueur \(6k\) par une matrice \(6\times6\).

Pour le canal de clôture :
\[
\begin{array}{c|c}
\text{profondeur}&\text{nouveau bloc}\\ \hline
6&010211\\
12&012222\\
18&010211\\
24&002111\\
30&110221
\end{array}
\]
et
\begin{equation}
 \boxed{
 w_{30}^{\pi}
 =010211|012222|010211|002111|110221.}
 \label{p12-83-c9-s1:eq:w30-pi-autonomo}
\end{equation}
La frontière suivante est
\[
 u_5^\pi=222220.
\]
Avec les autres canaux,
\[
 R_{36}=
 \begin{pmatrix}
 222220\\021222\\102011
 \end{pmatrix}.
\]
Le symbole \(R_{36}\) enregistre le sixième bloc de trois histoires et ouvre une
frontière ; ce n’est pas un terminal.

La branche certifiée de vingt blocs commence par
\begin{align*}
\pi:\;&010211|012222|010211|002111|110221|222220|111201|\\
&212121|200121|100100|101222|022212|012012|111210|\\
&121011|200220|120210|000101|022010|020111.
\end{align*}
C’est pourquoi ni \(w_{30}\) ni \(R_{36}\) ne terminent la généalogie.

\subsection{Décomposition orientée de la microfibre de \texorpdfstring{\(\pi\)}{pi} : région centrale, \(4\) régions latérales et multiplicité \texorpdfstring{\(432+4\cdot144=1008\)}{432+4x144=1008}}

La microfibre
\[
 \mathscr R_\pi^{\rm micro}
 :=\{U\in\mathcal U_6:\Psi(U)=010211\}
\]
contient exactement cinq régions :
\[
\begin{array}{c|c|c|c}
U_6&E_{\rm sig}&\operatorname{mult}&\text{rôle}\\ \hline
501|614|498|169|272|272&555555&432&\perp\\
810|923|870|169|272|272&339555&144&++\\
810|983|810|169|272|272&393555&144&++\\
870|923|810|169|272|272&933555&144&++\\
870|923|810|418|674|521&933717&144&--
\end{array}
\]
En conséquence,
\begin{equation}
 \boxed{
 5=1_\perp+3_{++}+1_{--},
 \qquad
 1008=432+4\cdot144.}
 \label{p12-83-c9-s1:eq:descomposiciones-cinco-regiones-pi}
\end{equation}
La première identité classe des fonctions d’orientation. La seconde compte
des multiplicités de provenance. Elles ne comptent pas cinq copies indiscernables.

La région transversale est
\[
 501|614|498|169|272|272,
 \qquad E_{\rm sig}=555555,
 \qquad\operatorname{mult}=432.
\]
Le retour muté est
\[
 870|923|810|418|674|521,
 \qquad E_{\rm sig}=933717,
 \qquad\operatorname{mult}=144.
\]
Les intervertir conserverait le recensement \(3/1/1\), mais fausserait la structure
régionale complète.

\subsection{Descente d'un bloc nonadique complet au multigraphe de transfert de phase}

Pour \(U_6=(U_1,\ldots,U_6)\), soit
\[
 A(U_6)=(U_1,U_2,U_3),
 \qquad
 B(U_6)=(U_4,U_5,U_6).
\]
Chaque émission définit une arête entre \(A(U_6)\) et \(B(U_6)\). Le multigraphe
brut possède \(96\) sommets, répartis en composantes de tailles
\[
 28,28,28,4,4,4.
\]
On passe au quotient de la phase interne par
\[
 (a,b,c)\sim(b,c,a)\sim(c,a,b).
\]
Le représentant est la rotation lexicographiquement minimale. Le graphe
quotient possède \(32\) sommets et deux composantes de tailles \(28\) et \(4\).

Les ancrages sont
\[
\begin{aligned}
U_{\rm APP}&=903|850|850|252|109|252,\\
U_e&=850|963|890|149|252|212,\\
U_\varphi&=830|943|830|109|252|252.
\end{aligned}
\]
Nous numérotons les régions
\[
\begin{aligned}
U_\pi^{(1)}&=810|923|870|169|272|272,\\
U_\pi^{(2)}&=810|983|810|169|272|272,\\
U_\pi^{(3)}&=870|923|810|169|272|272,\\
U_\pi^{(4)}&=870|923|810|418|674|521,\\
U_\pi^{(5)}&=501|614|498|169|272|272.
\end{aligned}
\]

\subsection{Les \(5\) cycles orientés de clôture}

\subsubsection{Clôture de \texorpdfstring{\(U_\pi^{(1)}\)}{U-pi 1}}
\begin{align*}
&252|109|252|903|850|850\to
903|850|850|252|109|252\to\\
&252|109|252|923|870|810\to
810|923|870|169|272|272\to\\
&272|169|272|963|890|850\to
850|963|890|149|252|212\to\\
&212|149|252|943|830|830\to
830|943|830|109|252|252\to
252|109|252|903|850|850.
\end{align*}

\subsubsection{Clôture de \texorpdfstring{\(U_\pi^{(2)}\)}{U-pi 2}}
\begin{align*}
&903|850|850|252|109|252\to
272|169|272|903|850|850\to\\
&272|169|272|983|810|810\to
810|983|810|169|272|272\to\\
&272|169|272|963|890|850\to
850|963|890|149|252|212\to\\
&212|149|252|943|830|830\to
830|943|830|109|252|252\to
903|850|850|252|109|252.
\end{align*}

\subsubsection{Clôture de \texorpdfstring{\(U_\pi^{(3)}\)}{U-pi 3}}
\begin{align*}
&903|850|850|252|109|252\to
292|189|232|903|850|850\to\\
&292|189|232|923|810|870\to
870|923|810|169|272|272\to\\
&272|169|272|963|890|850\to
850|963|890|149|252|212\to\\
&212|149|252|943|830|830\to
830|943|830|109|252|252\to
903|850|850|252|109|252.
\end{align*}

\subsubsection{Clôture de \texorpdfstring{\(U_\pi^{(4)}\)}{U-pi 4}}
\begin{align*}
&903|850|850|252|109|252\to
272|169|272|903|850|850\to\\
&272|169|272|963|890|850\to
850|963|890|149|252|212\to\\
&212|149|252|923|810|870\to
870|923|810|418|674|521\to\\
&943|830|830|521|418|674\to
830|943|830|109|252|252\to
903|850|850|252|109|252.
\end{align*}

\subsubsection{Clôture de \texorpdfstring{\(U_\pi^{(5)}\)}{U-pi 5}}
\begin{align*}
&252|109|252|903|850|850\to
903|850|850|252|109|252\to\\
&614|498|501|252|109|252\to
501|614|498|169|272|272\to\\
&272|169|272|963|890|850\to
850|963|890|149|252|212\to\\
&212|149|252|943|830|830\to
830|943|830|109|252|252\to
252|109|252|903|850|850.
\end{align*}

\begin{theorem}[Minimalité des cinq clôtures]
\label{p12-83-c9-s1:thm:minimalidad-cierres-pi-autonoma}
Pour chaque \(i\in\{1,\ldots,5\}\), la longueur minimale d’une marche fermée
qui parcourt simultanément
\[
 U_\pi^{(i)},\quad U_e,\quad U_\varphi,\quad U_{\rm APP}
\]
dans le multigraphe quotient est huit.
\end{theorem}

\begin{proof}
Pour \(i\) fixé, soit
\[
 E_*=\{U_\pi^{(i)},U_e,U_\varphi,U_{\rm APP}\}.
\]
On considère l’espace fini
\[
 \mathscr S_i=\{(v,M):v\in V(G),\ M\subseteq E_*\}.
\]
La première composante enregistre le sommet et la seconde les arêtes cibles
déjà parcourues. Si une arête \(e\) mène de \(v\) à \(v'\), la transition est
\[
 (v,M)\longmapsto
 \bigl(v',M\cup(E_*\cap\{e\})\bigr).
\]
La recherche en largeur énumère les longueurs par ordre croissant. Le premier
retour à \((v,E_*)\) apparaît à la longueur huit pour les cinq régions.
Les marches imprimées réalisent la borne ; l’optimalité de la recherche exclut
une longueur inférieure.
\end{proof}

La fonction de clôture est encodée par
\[
 \mathfrak C_\pi=
 \bigl(\mathscr R_\pi^{\rm micro},
 w_\pi,\rho_\pi,m_\pi,\ell_{\rm cl}\bigr),
\]
\[
 \rho_\pi=(\perp,++ ,++ ,++ ,--),
 \qquad
 m_\pi=(432,144,144,144,144),
 \qquad
 \ell_{\rm cl}=8.
\]

\subsection{Opérateur d’incidence et demi-tour orienté}

L’ombre d’incidence est
\[
 A_W=
 \begin{pmatrix}
 0&1&1&1&1&1\\
 1&0&1&2&2&1\\
 1&1&0&1&2&2\\
 1&2&1&0&1&2\\
 1&2&2&1&0&1\\
 1&1&2&2&1&0
 \end{pmatrix}.
\]
La multiplication dans \(\mathbb F_3\) donne
\begin{equation}
 \boxed{A_W^2=2I_6=-I_6.}
 \label{p12-83-c9-s1:eq:aw-cuadrado-menos-identidad-pi}
\end{equation}
Les produits entre lignes et colonnes distinctes s’annulent modulo trois,
tandis que chaque produit diagonal vaut deux. Une application réalise un quart
de tour ternaire ; deux applications réalisent l’inversion orientée.

\subsection{La pente 1/5 imposée par les \(5\) régions}

Soit
\[
 \lambda=(\lambda_1,\ldots,\lambda_5)^{\mathsf T}
 \in\mathbb Q_{\geq0}^5,
 \qquad
 \mathbf1^{\mathsf T}\lambda=1.
\]
Oublier l’étiquette régionale applique le symétriseur
\[
 P_5=\frac15\mathbf1\mathbf1^{\mathsf T}.
\]
Pour tout vecteur normalisé,
\[
 P_5\lambda=\frac15\mathbf1.
\]
L’invariance \(P_5\lambda=\lambda\) a donc une unique solution :
\begin{equation}
 \boxed{\lambda=\frac15\mathbf1,\qquad t_1=\frac15.}
 \label{p12-83-c9-s1:eq:pendiente-primitiva-pi}
\end{equation}
L’entier cinq provient de la microfibre régionale ; il n’est pas reconnu dans une
formule de \(\pi\).

\subsection{Composition rationnelle et compensateur \texorpdfstring{\(239\)}{239}}

La loi de composition des pentes est
\[
 t\star s:=\frac{t+s}{1-ts}.
\]
Le premier doublement donne
\[
 t_2=t_1\star t_1
 =\frac{2/5}{1-1/25}
 =\frac5{12}.
\]
Le second donne
\[
 t_4=t_2\star t_2
 =\frac{10/12}{1-25/144}
 =\frac{120}{119}.
\]
Le retour à la diagonale unitaire exige \(q>0\) tel que
\[
 \frac{120/119-1/q}{1+(120/119)/q}=1.
\]
La multiplication par \(119q\) produit
\[
 120q-119=119q+120,
\]
et impose
\begin{equation}
 \boxed{q=239.}
 \label{p12-83-c9-s1:eq:compensador-239-pi}
\end{equation}

\subsection{Identité gaussienne et 1ʳᵉ demi-révolution}

On définit
\begin{equation}
 \Theta_\pi
 :=16\arctan\frac15-4\arctan\frac1{239}.
 \label{p12-83-c9-s1:eq:theta-pi-autonoma}
\end{equation}
Les carrés successifs sont
\[
\begin{array}{c|r@{}l}
\text{puissance}&\multicolumn{2}{c}{\text{entier gaussien}}\\ \hline
(5+i)^2&24&+10i\\
(5+i)^4&476&+480i\\
(5+i)^8&-3824&+456960i\\
(5+i)^{16}&-208797818624&-3494830080i\\
(239-i)^2&57120&-478i\\
(239-i)^4&3262465916&-54606720i
\end{array}
\]
et la multiplication finale donne
\begin{equation}
 \boxed{
 (5+i)^{16}(239-i)^4
 =-681386607803576157184.}
 \label{p12-83-c9-s1:eq:identidad-gaussiana-pi-autonoma}
\end{equation}
La partie imaginaire est nulle et la partie réelle est négative. Par conséquent,
\[
 \Theta_\pi\equiv\pi\pmod{2\pi}.
\]
Para \(0<x<1\),
\[
 x-\frac{x^3}{3}<\arctan x<x.
\]
D’où
\[
 \Theta_\pi>
 16\left(\frac15-\frac1{3\cdot5^3}\right)-\frac4{239}>3
\]
et
\[
 \Theta_\pi<\frac{16}{5}<4.
\]
L’unique orientation négative dans cet intervalle est le premier demi-tour
positif :
\begin{equation}
 \boxed{\Theta_\pi=\pi.}
 \label{p12-83-c9-s1:eq:reconocimiento-pi-exacto}
\end{equation}

L’ordre de la reconnaissance archimédienne ultérieure est
\[
 \text{cinq régions}
 \to\frac15
 \to\frac5{12}
 \to\frac{120}{119}
 \to239
 \to\Theta_\pi
 \to\pi.
\]

\subsection{Encadrements rationnels de précision arbitraire}

Pour \(0<x<1\), soit
\[
 S_n(x)=\sum_{r=0}^{n}(-1)^r\frac{x^{2r+1}}{2r+1}.
\]
Le critère de Leibniz donne
\[
 S_{2m+1}(x)<\arctan x<S_{2m}(x).
\]
On définit
\[
 a_m^-=S_{2m+1}(1/5),\quad a_m^+=S_{2m}(1/5),
\]
\[
 b_m^-=S_{2m+1}(1/239),\quad b_m^+=S_{2m}(1/239),
\]
et
\begin{equation}
 \ell_{\pi,m}=16a_m^- -4b_m^+,
 \qquad
 h_{\pi,m}=16a_m^+ -4b_m^-.
 \label{p12-83-c9-s1:eq:horquilla-pi-autonoma}
\end{equation}
Alors
\[
 \ell_{\pi,m}<\pi<h_{\pi,m}
\]
et
\begin{equation}
 h_{\pi,m}-\ell_{\pi,m}
 =
 \frac{16}{(4m+3)5^{4m+3}}
 +\frac{4}{(4m+3)239^{4m+3}}
 \longrightarrow0.
 \label{p12-83-c9-s1:eq:anchura-horquilla-pi}
\end{equation}
Les bornes inférieures croissent et les bornes supérieures décroissent.

Comme \(3<\pi<4\), le lecteur ternaire agit sur
\[
 r_\pi:=\pi-3\in(0,1),
\]
avec encadrement
\[
 [\ell_{\pi,m}-3,h_{\pi,m}-3].
\]

\subsection{Certification archimédienne ultérieure de la fibre déjà engendrée}

Supposons que la relation interne \(\operatorname{Ext}\) et la survie
coinductive aient déjà produit, sans utiliser Machin ni un développement tabulé,
l’extension compatible \(N_{K+1}\) du préfixe \(N_K\), de longueur
\(L_K=6K\). Les \(729\) sous-cylindres de sa fibre ambiante sont
\[
 I_{K,u}=
 \left[
 \frac{729N_K+u}{3^{L_K+6}},
 \frac{729N_K+u+1}{3^{L_K+6}}
 \right),
 \qquad0\leq u<729.
\]
L’ordre \(m=m(K)\) est ensuite augmenté jusqu’à ce que l’encadrement de Leibniz
localise le sous-cylindre déjà produit. Soient
\[
 \ell_{\pi,K}:=\ell_{\pi,m(K)}-3,
 \qquad
 h_{\pi,K}:=h_{\pi,m(K)}-3.
\]
Les extrémités entières compatibles sont
\[
 j_{\min}=
 \max\!\left(
 729N_K,\left\lfloor
 \ell_{\pi,K}3^{L_K+6}\right\rfloor
 \right),
\]
\[
 j_{\max}=
 \min\!\left(
 729N_K+728,\left\lceil
 h_{\pi,K}3^{L_K+6}\right\rceil-1
 \right).
\]
La condition de localisation unitaire est
\begin{equation}
 \boxed{j_{\min}=j_{\max},\qquad N_{K+1}=j_{\min}.}
 \label{p12-83-c9-s1:eq:seleccion-unitaria-subcilindro-pi}
\end{equation}
La première égalité certifie la localisation unitaire ; la seconde vérifie
que cette localisation coïncide avec l’extension engendrée intérieurement. Aucune
des deux ne définit la transition TPK ni ne sélectionne parmi les \(729\) états.
L’existence et l’unicité procèdent de \(\operatorname{Ext}\) et de la
survie ; Machin et les encadrements de Leibniz sont des reconnaissances
ultérieures.

\subsection{Prolongement coinductif de la section de \(\pi\) par le système inverse des histoires survivantes}

Soit \(\mathcal H_m^\pi\) l’ensemble des histoires enrichies à la profondeur
\(m\), et soit
\[
 p_{n,m}:\mathcal H_n^\pi\longrightarrow\mathcal H_m^\pi
\]
la troncature. Nous définissons
\[
 \operatorname{Surv}_m^{(N),\pi}
 :=p_{N,m}(\mathcal H_N^\pi),
\]
\[
 \operatorname{Surv}_m^\pi
 :=\bigcap_{N\geq m}\operatorname{Surv}_m^{(N),\pi}.
\]
Lorsque la relation \(\operatorname{Ext}\) est non vide, univoque et
compatible à tous les niveaux, les ensembles sont des singletons et satisfont
\[
 p_{m+1,m}(\operatorname{Surv}_{m+1}^{\pi})
 =\operatorname{Surv}_{m}^{\pi}.
\]
Dans ce domaine de prolongeabilité, il existe donc une unique section
\[
 x_\pi\in
 X_\pi:=\varprojlim_m\operatorname{Surv}_m^\pi.
\]
Les cylindres sont emboîtés et leurs diamètres sont \(3^{-6K}\), qui tend vers
zéro. Leur intersection contient un unique point, identifié par le caractère
angulaire à \(r_\pi=\pi-3\).

Pour toute profondeur décimale \(M\), on peut choisir un niveau \(K(M)\) dont le
cylindre est contenu dans une seule cellule décimale de diamètre \(10^{-M}\).
Les préfixes profonds se tronquent exactement aux précédents.

\subsection{Caractérisation du canal coinductif de \(\pi\) : sélection de \(w_6^\pi\), \(5\) régions, inversion orientée et unicité de la section limite}

\begin{theorem}[Caractérisation du canal de clôture]
\label{p12-83-c9-s1:thm:caracterizacion-final-pi-autonoma}
À partir du catalogue, de l’action \(D_3\), de la jauge orientée, des
opérateurs \(L_0,L_1,A_W\), de la loi des pentes et de la filtration rationnelle :
\begin{enumerate}
 \item on sélectionne \(w_6^\pi=010211\) ;
 \item on construit cinq régions avec les identités de
       \eqref{p12-83-c9-s1:eq:descomposiciones-cinco-regiones-pi} ;
 \item chaque région appartient à une clôture minimale de longueur huit ;
 \item \(A_W^2=-I_6\) fixe l’inversion orientée ;
 \item la symétrisation régionale force \(1/5\) ;
 \item la composition force \(120/119\) ;
 \item le retour unitaire force \(239\) ;
 \item l’identité gaussienne reconnaît le premier demi-tour ;
 \item les encadrements et la fibre de \(729\) éléments produisent une section
       compatible à toute profondeur finie.
\end{enumerate}
Son évaluation est
\[
 \boxed{\pi=16\arctan\frac15-4\arctan\frac1{239}.}
\]
\end{theorem}

La chaîne complète est
\[
 \mathrm{APP}\to\mathrm{TRIT}\to\mathrm{TPK}
 \to104\,976\to468\to243\to w_6^\pi
 \to w_{30}^\pi\to R_{36}\to\Gamma_9
 \to X_\pi\to\pi.
\]

\begin{criterion}[Réfutation du canal de clôture]
La dérivation est réfutée si :
\begin{enumerate}
 \item l’orbite de clôture cesse d’être unique ;
 \item l’une des cinq régions disparaît ou son identité complète change ;
 \item l’une des cinq clôtures a une longueur inférieure à huit ;
 \item \(A_W^2\ne2I_6\);
 \item le symétriseur admet une autre distribution normalisée ;
 \item les compositions ne produisent pas \(5/12\), \(120/119\) et \(239\) ;
 \item l’identité gaussienne échoue ;
 \item un encadrement cesse de contenir le caractère angulaire ;
 \item un niveau ne parcourt pas les \(729\) éléments de la fibre ambiante ou ne
       localise pas exactement une extension ;
 \item un développement cible intervient avant la reconnaissance.
\end{enumerate}
\end{criterion}
\fi

La construction complète du caractère de clôture, des cinq régions et de la
biographie structurelle de \(\pi\) a été établie dans le
\cref{chap:83-08}. À partir de ce même caractère déjà construit, le présent
chapitre prolonge maintenant le raffinement vers \(e\) et \(\varphi\), sans
réinitialiser la généalogie ni répéter sa première étape.


% Unidad U013. Biografía estructural completa del carácter exponencial.

\section{Propagation, retour avec mémoire et construction du caractère exponentiel}
\label{p12-83-c9-s2:chap:biografia-completa-e}

\subsection{Sélecteur orbital fini : existence et unicité de l’émission compatible}

Soit \(\mathcal W=\mathbb F_3^6\). Pour une émission
\(U=(u_1,\ldots,u_6)\), le mot visible est
\[
 \Psi(U)=((-u_1)\bmod3,\ldots,(-u_6)\bmod3).
\]
Sur \(\mathcal W\), l’action diédrale est donnée par
\[
\begin{aligned}
 r(u_1,u_2,u_3,u_4,u_5,u_6)
 &=(u_3,u_1,u_2,u_5,u_6,u_4),\\
 s(u_1,u_2,u_3,u_4,u_5,u_6)
 &=(u_4,u_5,u_6,u_1,u_2,u_3).
\end{aligned}
\]
Pour \(w\in\mathcal W\), nous définissons
\[
 n(w)=\#\Psi^{-1}(w),
 \qquad
 M(w)=\sum_{U\in\Psi^{-1}(w)}\operatorname{mult}(U),
\]
et
\[
 \nu(w)=
 \bigl(
 \#\{j:w_j=0\},
 \#\{j:w_j=1\},
 \#\{j:w_j=2\}
 \bigr).
\]

\begin{theorem}[Sélection structurelle du canal exponentiel]
\label{p12-83-c9-s2:thm:seleccion-estructural-e}
Parmi les \(43\) orbites du catalogue visible, il existe une unique orbite dont les
fibres satisfont simultanément :
\[
 n(w)=1,\qquad M(w)=144,\qquad
 \mathfrak D(w)=\{0,3,6\},\qquad
 \nu(w)=(2,3,1).
\]
Le calibre
\[
 \operatorname{diag}_c=6,
 \qquad \operatorname{card}=ES
\]
y sélectionne
\[
 \boxed{w_6^e=201101.}
\]
\end{theorem}

\begin{proof}
Les quatre conditions sont évaluées sur les \(243\) mots du catalogue,
non sur un échantillon. Une unique orbite satisfait le prédicat conjoint. Dans
cette orbite, une seule émission possède à la fois la diagonale \(6\) et l’orientation
\(ES\). L’énumération utilise uniquement des champs du catalogue APP--TRIT et ne
consulte pas de développement de \(e\).
\end{proof}

\subsection{Relèvement visible et chaîne de blocs}

La récurrence \(L_0,L_0,L_1,L_1\) produit
\[
 w_{30}^{e}
 =201101|121221|102011|012222|102011,
\]
et la frontière conjointe apporte
\[
 u_5^e=021222.
\]
Le registre de trente-six trits devient
\[
 w_{36}^{e}
 =201101|121221|102011|012222|102011|021222.
\]

\subsection{Signature décimale et carré positionnel}

La lecture initiale de douze chiffres est
\[
 \mathbf d_e^{(12)}
 =718281828459
 =7\,\underbrace{1828\,1828}_{\text{carré local}}\,459.
\]
L’égalité centrale est la concaténation
\[
 1828\mathbin{\|}1828,
\]
non un produit arithmétique ni le début supposé d’une période.

\begin{definition}[Carré positionnel]
Un mot \(d_1\cdots d_{12}\) contient un carré de longueur \(\ell\)
à la position \(p\) lorsque
\[
 d_p\cdots d_{p+\ell-1}
 =d_{p+\ell}\cdots d_{p+2\ell-1}.
\]
La position et les contextes latéraux appartiennent à la signature.
\end{definition}

\begin{table}[htbp]
\centering
\caption{Recensement exhaustif des carrés dans les lectures de douze chiffres.}
\label{p12-83-c9-s2:tab:censo-cuadrados-e}
\begin{tabular}{c*{6}{r}}
\toprule
Longueur \(\ell\)&1&2&3&4&5&6\\
\midrule
Occurrences&240&21&3&2&0&0\\
Mots qui les contiennent&153&19&3&2&0&0\\
\bottomrule
\end{tabular}
\end{table}

L’autre mot contenant un carré de longueur quatre est
\[
 575157518472
 =\underbrace{5751\,5751}_{\text{position initiale }1}\,8472.
\]

\begin{proposition}[Unicité positionnelle]
\label{p12-83-c9-s2:prop:cuadrado-posicional-e}
La lecture du canal \(e\) est la seule qui possède un carré de longueur
quatre après un préfixe de longueur un et avec les contextes \(7\) et \(459\).
\end{proposition}

\begin{proof}
Le recensement ne laisse que deux mots contenant des carrés de longueur quatre. L’un
commence son carré à la première position ; l’autre après le préfixe \(7\). La
position et les contextes séparent les deux possibilités.
\end{proof}

\subsection{Retour visible et rupture de mémoire}

Dans le relèvement
\[
 201101|121221|102011|012222|102011
\]
les troisième et cinquième blocs coïncident :
\[
 B_3=B_5=102011.
\]
Leurs préétats modulo \(27\) sont cependant
\[
 p_3=(25,11,7,17,8,16),
 \qquad
 p_5=(7,8,4,23,26,7),
\]
et les quotients ultérieurs :
\[
 q_3=(6,13,9,2,13,1),
 \qquad
 q_5=(15,0,3,19,23,25).
\]
On vérifie \(p_3\ne p_5\) et \(q_3\ne q_5\).

\begin{figure}[htbp]
\centering
\begin{tikzpicture}[
 node distance=11mm and 7mm,
 every node/.style={font=\small},
 block/.style={draw,rounded corners,minimum width=17mm,minimum height=8mm},
 state/.style={draw,dashed,rounded corners,align=center,inner sep=3pt},
 >=stealth
]
\node[block] (b1) {\(201101\)};
\node[block,right=of b1] (b2) {\(121221\)};
\node[block,right=of b2] (b3) {\(102011\)};
\node[block,right=of b3] (b4) {\(012222\)};
\node[block,right=of b4] (b5) {\(102011\)};
\draw[->] (b1)--(b2);
\draw[->] (b2)--(b3);
\draw[->] (b3)--(b4);
\draw[->] (b4)--(b5);
\node[state,below=of b3] (p3)
 {\(p_3=(25,11,7,17,8,16)\)\\
  \(q_3=(6,13,9,2,13,1)\)};
\node[state,below=of b5] (p5)
 {\(p_5=(7,8,4,23,26,7)\)\\
  \(q_5=(15,0,3,19,23,25)\)};
\draw[->] (p3)--(b3);
\draw[->] (p5)--(b5);
\draw[<->,thick] (b3) to[bend left=28]
 node[above]{même mot visible} (b5);
\end{tikzpicture}
\caption{Retour visible de \(102011\) avec renouvellement du préétat et du
quotient.}
\label{p12-83-c9-s2:fig:ruptura-memoria-e}
\end{figure}

\begin{theorem}[Retour visible sans périodicité de l’état]
\label{p12-83-c9-s2:thm:retorno-visible-no-periodico-e}
\(B_3=B_5\) est un retour de la projection visible, non une orbite périodique
de l’état enrichi.
\end{theorem}

\begin{proof}
Le retour de l’état complet exigerait l’égalité de toutes ses coordonnées.
Les inégalités du préétat et du quotient l’excluent. De même,
\(1828|1828\) ne commence pas une période décimale : après le deuxième bloc apparaît
\(4\), tandis qu’une continuation périodique exigerait \(1\).
\end{proof}

\subsection{Certification factorielle ultérieure de la section exponentielle coinductive}

Soit \(\mathcal A\) une algèbre commutative unitaire sur \(\mathbb Q\), et
soit \(P_t\in\mathcal A[[t]]\) le propagateur formel. La concaténation
temporelle impose
\[
 P_{s+t}=P_sP_t,
 \qquad
 P_0=1.
\]
Écrivons
\[
 P_t=\sum_{n=0}^{\infty}a_nt^n.
\]
L’unité fixe \(a_0=1\), et le transport élémentaire normalisé fixe
\(a_1=1\).

\begin{theorem}[Récurrence factorielle]
\label{p12-83-c9-s2:thm:recurrencia-factorial-e}
\[
 (n+1)a_{n+1}=a_n
 \qquad(n\geq0),
\]
et, par conséquent,
\[
 a_n=\frac1{n!},
 \qquad
 P_t=\sum_{n=0}^{\infty}\frac{t^n}{n!}.
\]
En \(t=1\),
\[
 \boxed{P_1=e.}
\]
\end{theorem}

\begin{proof}
Le coefficient de \(s^nt\) dans \(P_{s+t}\) est
\[
 \binom{n+1}{n}a_{n+1}=(n+1)a_{n+1}.
\]
Dans \(P_sP_t\), il est \(a_na_1=a_n\). L’égalité des séries produit la
récurrence. Une récurrence à partir de \(a_0=1\) donne \(a_n=1/n!\).
\end{proof}

\begin{corollary}[Équation différentielle]
\[
 \frac{d}{dt}P_t=P_t,
 \qquad P_0=1.
\]
\end{corollary}

\begin{proof}
\[
 P_t'=\sum_{n\geq0}(n+1)a_{n+1}t^n
 =\sum_{n\geq0}a_nt^n=P_t.
\]
\end{proof}

\subsection{Emboîtement cofinal des encadrements rationnels et convergence vers la valeur publiée}

Soit
\[
 S_N=\sum_{n=0}^{N}\frac1{n!}.
\]

\begin{theorem}[Borne factorielle explicite]
\label{p12-83-c9-s2:thm:cota-factorial-e}
Pour \(N\geq1\),
\begin{equation}
 \boxed{
 S_N<e<
 S_N+\frac{N+2}{(N+1)(N+1)!}.}
 \label{p12-83-c9-s2:eq:horquilla-factorial-e}
\end{equation}
\end{theorem}

\begin{proof}
Le reste est
\[
 R_N=\sum_{k=0}^{\infty}\frac1{(N+1+k)!}>0.
\]
Para \(k\geq0\),
\[
 (N+1+k)!
 \geq(N+1)!(N+2)^k.
\]
Par conséquent,
\[
 R_N
 \leq\frac1{(N+1)!}
 \sum_{k=0}^{\infty}\frac1{(N+2)^k}
 =\frac{N+2}{(N+1)(N+1)!}.
\]
L’inégalité est stricte à partir de \(k=2\).
\end{proof}

On définit
\[
 J_N^{(e)}=
 \left[
 S_N,\,
 S_N+\frac{N+2}{(N+1)(N+1)!}
 \right]
\]
et, pour la partie fractionnaire,
\[
 \widetilde J_N^{(e)}=J_N^{(e)}-2.
\]
La largeur converge vers zéro. Pour chaque \(M\), il existe \(N\) tel que
\[
 \operatorname{diam}J_N^{(e)}<10^{-M}.
\]
Si les extrémités appartiennent à une même cellule décimale, les \(M\) premiers
chiffres sont certifiés par arithmétique rationnelle.

\subsection{Prolongement coinductif de la section exponentielle par les sous-cylindres ternaires du Topological Phase Kernel}

Au niveau \(K\), TPK et la relation \(\operatorname{Ext}\) ont déjà engendré
un unique sous-cylindre compatible. On prend le plus petit \(N=N_e(K)\) pour lequel
l’encadrement factoriel est contenu dans ce sous-cylindre. \(N_e(K)\) est un
ordre de certification analytique, non un indice générateur ni une règle de
sélection de la fibre.

Les cylindres TPK sont compatibles par troncature indépendamment du
semi-groupe. Le théorème du semi-groupe identifie ensuite leur évaluation
archimédienne à \(e-2\) ; la partie entière retrouve \(e\). Si l’orbite,
\(\operatorname{Ext}\) ou la troncature échouent, la génération échoue ; si la
récurrence ou la borne factorielle échouent, c’est ce certificat analytique qui échoue, non la
généalogie APP--TRIT--TPK.

\subsection{Certificat adversarial de la génération de \(e\) : unicité orbitale, récurrence factorielle, encadrements rationnels et prolongation cylindrique}

\begin{criterion}[Réfutation du canal exponentiel]
La construction est réfutée si :
\begin{enumerate}
 \item une autre orbite satisfait le prédicat fini complet ;
 \item la signature \(7[1828\,1828]459\) n’est pas positionnellement unique ;
 \item les préétats ou les quotients des deux blocs \(102011\) coïncident ;
 \item le semi-groupe ne force pas \((n+1)a_{n+1}=a_n\) ;
 \item la borne \eqref{p12-83-c9-s2:eq:horquilla-factorial-e} échoue ;
 \item un niveau ne localise pas de manière unitaire une extension compatible ;
 \item un développement tabulé intervient avant la reconnaissance.
\end{enumerate}
\end{criterion}


% Unidad U014. Biografía estructural completa del carácter de autoescala.

\section{Incidence modale, demi-droite de Perron et construction du caractère d'auto-échelle}
\label{p12-83-c9-s3:chap:biografia-completa-phi}

\subsection{Sélection finie du germe d’auto-échelle}

Soit \(\mathcal U_6\subset\{0,\ldots,999\}^6\) le catalogue des
\(468\) émissions distinctes construit dans le
\cref{chap:semillas-emisor-54}, et soit
\[
 \Psi:\mathcal U_6\longrightarrow\mathbb F_3^6,
 \qquad
 \Psi(U)=((-U_j)\bmod3)_{j=1}^{6}.
\]
Pour \(w\in\Psi(\mathcal U_6)\), rappelons les invariants
\[
 n(w)=\#\Psi^{-1}(w),
 \qquad
 M(w)=\sum_{U\in\Psi^{-1}(w)}\operatorname{mult}(U),
\]
et le vecteur d’occupation
\[
 \nu(w)=
 \bigl(\#\{j:w_j=0\},\#\{j:w_j=1\},\#\{j:w_j=2\}\bigr).
\]

\begin{theorem}[Sélection orbitale de l’auto-échelle]
\label{p12-83-c9-s3:thm:seleccion-orbital-phi}
L’action diédrale sur \(\Psi(\mathcal U_6)\) contient une unique orbite
qui satisfait le prédicat d’auto-échelle
\[
 n(w)=1,
 \qquad
 M(w)=144,
 \qquad
 \nu(w)=(2,2,2).
\]
La jauge interne
\[
 \operatorname{diag}_c=6,
 \qquad
 \operatorname{card}=ES
\]
sélectionne dans cette orbite le mot
\[
 \boxed{w_6^{\varphi}=121200.}
\]
Son unique émission est
\[
 \boxed{U_{\varphi}=(830,943,830,109,252,252),}
\]
de multiplicité \(144\), de signature multiplicative \(555933\), d’orientation
additive \(ES\), de support multiplicatif de cardinal six et de type résiduel
\(A\).
\end{theorem}

\begin{proof}
La projection directe donne
\[
 U_{\varphi}\bmod3=(2,1,2,1,0,0),
 \qquad
 -U_{\varphi}\bmod3=(1,2,1,2,0,0).
\]
L’énumération est effectuée sur les \(243\) mots visibles et leurs
\(43\) orbites, non sur une sélection d’exemples. Le prédicat conjoint
\(n=1\), \(M=144\), \(\nu=(2,2,2)\) laisse une orbite. En son sein, la
diagonale et l’orientation déclarées laissent une émission. À cette étape,
ni l’équation \(x^2-x-1=0\), ni une table décimale, ni le nom
conventionnel de la constante n’ont été introduits.
\end{proof}

\subsection{Relèvement par blocs et 1ʳᵉ frontière coinductive}

Soit \(b_1=w_6^{\varphi}\), et soit
\(\varepsilon=(0,0,1,1)\). Avec les matrices \(L_0,L_1\) construites dans le
\cref{hmtctx:p12:chap:orbitas-elevacion-r36}, nous définissons
\[
 b_{k+1}=b_kL_{\varepsilon_k}\pmod3
 \qquad(1\leq k\leq4),
\]
et concaténons
\[
 w_{6(k+1)}^{\varphi}=w_{6k}^{\varphi}\mathbin{\|}b_{k+1}.
\]
Le calcul matriciel produit, bloc par bloc,
\[
\begin{aligned}
 b_1&=121200,\\
 b_2&=112202,\\
 b_3&=121020,\\
 b_4&=010210,\\
 b_5&=010200,
\end{aligned}
\]
et, par conséquent,
\begin{equation}
 \boxed{
 w_{30}^{\varphi}
 =121200|112202|121020|010210|010200.}
 \label{p12-83-c9-s3:eq:w30-phi}
\end{equation}
La neutralité duale de la frontière ajoute
\[
 u_5^{\varphi}=102011,
\]
de sorte que
\begin{equation}
 \boxed{
 w_{36}^{\varphi}
 =121200|112202|121020|010210|010200|102011.}
 \label{p12-83-c9-s3:eq:w36-phi}
\end{equation}

\begin{remark}[Deux usages distincts d’un mot visible]
Le mot \(102011\) est également visible dans la biographie de \(e\). L’égalité
d’une projection n’identifie pas les états enrichis : le canal,
le préfixe, le quotient, la retenue, la phase et la mémoire appartiennent au type
de l’objet. Cette séparation est obligatoire avant de construire tout
caractère réel.
\end{remark}

\subsection{Du calendrier tritique au multigraphe d’incidence}

\subsubsection{Incidence modale primitive \(F_{\rm av}\) et dynamique de mémoire d'ordre \(1\)}

Le calendrier tritique primitif est
\[
 (+1,-1,0).
\]
Nous fixons deux modes visibles, \(\Sigma\) et \(\Pi\), et la règle
\begin{equation}
 +1:m\longmapsto\Sigma,
 \qquad
 -1:m\longmapsto\Pi,
 \qquad
 0:m\longmapsto m.
 \label{p12-83-c9-s3:eq:regla-modal-primitiva-phi}
\end{equation}
Avec la condition de clôture cyclique, l’orbite modale produite par
\eqref{p12-83-c9-s3:eq:regla-modal-primitiva-phi} est
\[
 \boxed{(\Sigma,\Pi,\Pi).}
\]
Les trois paires consécutives sont
\[
 \Sigma\longrightarrow\Pi,
 \qquad
 \Pi\longrightarrow\Pi,
 \qquad
 \Pi\longrightarrow\Sigma.
\]

\begin{definition}[Matrice d’incidence chronologique]
Pour \(i,j\in\{\Sigma,\Pi\}\), soit
\[
 N_3(i,j)=
 \#\{r\in\mathbb Z/3\mathbb Z:(m_r,m_{r+1})=(i,j)\}.
\]
L’indice \(3\) compte les instants du calendrier primitif ; il ne désigne pas une
puissance matricielle.
\end{definition}

\begin{theorem}[Descente nécessaire de l’incidence modale]
\label{p12-83-c9-s3:thm:descenso-incidencia-modal-phi}
Dans l’ordre \((\Sigma,\Pi)\),
\begin{equation}
 \boxed{
 N_3=
 \begin{pmatrix}
 0&1\\
 1&1
 \end{pmatrix}.}
 \label{p12-83-c9-s3:eq:N3-phi}
\end{equation}
Par répétition exacte du calendrier,
\begin{equation}
 \boxed{
 N_9=3N_3,
 \qquad
 N_{27}=9N_3,
 \qquad
 N_{108}=36N_3.}
 \label{p12-83-c9-s3:eq:incidencias-9-27-108-phi}
\end{equation}
Ces trois égalités sont des égalités de dénombrements d’arêtes ; elles n’affirment
ni \(N_9=N_3^3\), ni \(N_{27}=N_3^9\), ni \(N_{108}=N_3^{36}\).
\end{theorem}

\begin{proof}
\(\Sigma\to\Sigma\) n’apparaît pas ; chacune des trois autres paires apparaît
une fois. Cela détermine les quatre entrées de \(N_3\). Les calendriers de
longueurs \(9\), \(27\) et \(108\) contiennent respectivement \(3\), \(9\)
et \(36\) copies du bloc primitif. Chaque répétition multiplie tous les
dénombrements par le même entier, ce qui prouve
\eqref{p12-83-c9-s3:eq:incidencias-9-27-108-phi}.
\end{proof}

\subsubsection{Publication surfacique–volumétrique et rapport stable \(\varphi^{-2}\) sous la mesure de Parry}

L’étiquette modale et le feuillet géométrique sont des objets distincts. La carte
de publication est
\[
 \mathfrak g_{\rm av}(\Sigma)=\mathscr H_{\rm ar},
 \qquad
 \mathfrak g_{\rm av}(\Pi)=\mathscr H_{\rm vol}.
\]
En transportant \eqref{p12-83-c9-s3:eq:N3-phi} par \(\mathfrak g_{\rm av}\), nous obtenons
l’opérateur d’incidence
\begin{equation}
 \boxed{
 F_{\rm av}=
 \begin{pmatrix}
 0&1\\
 1&1
 \end{pmatrix}.}
 \label{p12-83-c9-s3:eq:Fav-derivada-phi}
\end{equation}
Ainsi, les deux transitions croisées, l’unique autorétention volumétrique et
l’absence d’autorétention aréale découlent du calendrier ; ce ne sont pas quatre
entrées indépendantes.

\begin{remark}[Portée de la descente]
L’oubli de la phase élémentaire ne produit pas, à lui seul, un quotient de Markov
fort du Topological Phase Kernel, parce que le mode visible ne détermine pas
lequel des trois opérateurs internes agira au pas suivant. Ce qui
descend en revanche sans hypothèse supplémentaire est le multigraphe des arêtes d’un
bloc complet. Son enveloppe minimale de mémoire un a pour matrice
d’adjacence \(F_{\rm av}\).
\end{remark}

\subsection{Demi-droite positive invariante et unicité de l’auto-échelle}

\begin{theorem}[Demi-droite de Perron dérivée]
\label{p12-83-c9-s3:thm:rayo-perron-phi}
Le cône positif de \(F_{\rm av}\) contient une unique demi-droite propre. Si nous la
normalisons comme \(v=(1,x)^{\mathsf T}\), alors
\[
 x^2-x-1=0,
 \qquad
 x>0,
\]
et, par conséquent,
\begin{equation}
 \boxed{x=\frac{1+\sqrt5}{2}=\varphi.}
 \label{p12-83-c9-s3:eq:phi-perron}
\end{equation}
\end{theorem}

\begin{proof}
L’équation \(F_{\rm av}v=\lambda v\) donne, coordonnée par coordonnée,
\[
 x=\lambda,
 \qquad
 1+x=\lambda x.
\]
En éliminant \(\lambda\), on obtient \(x^2-x-1=0\). Ses racines sont
\((1\pm\sqrt5)/2\), et seule la positive appartient à l’intérieur du cône. La
matrice \(F_{\rm av}\) est primitive, puisque \(F_{\rm av}^2\) a toutes ses
entrées positives ; le théorème de Perron--Frobenius garantit que cette demi-droite
positive est simple et unique. Le nom \(\varphi\) est attribué après cette
construction.
\end{proof}

\begin{corollary}[Équation scalaire d’auto-échelle]
\[
 \varphi=1+\frac1\varphi,
 \qquad
 \varphi^2=\varphi+1,
 \qquad
 \varphi^{-1}=\varphi-1.
\]
Ces identités sont des conséquences de \eqref{p12-83-c9-s3:eq:phi-perron} ; elles ne sélectionnent
ni le germe ni la matrice.
\end{corollary}

\subsection{Récurrence de Fibonacci de la demi-droite de Perron et emboîtement diophantien des encadrements rationnels}

Soit \(F_0=0\), \(F_1=1\) et
\[
 F_{n+1}=F_n+F_{n-1}\qquad(n\geq1).
\]

\begin{theorem}[Puissances exactes de l’incidence]
\label{p12-83-c9-s3:thm:potencias-fibonacci-phi}
Pour tout \(n\geq1\),
\begin{equation}
 \boxed{
 F_{\rm av}^{\,n}=
 \begin{pmatrix}
 F_{n-1}&F_n\\
 F_n&F_{n+1}
 \end{pmatrix}.}
 \label{p12-83-c9-s3:eq:potencias-Fav-fibonacci}
\end{equation}
\end{theorem}

\begin{proof}
Pour \(n=1\), l’identité est la définition de \(F_{\rm av}\). Si elle vaut
pour \(n\), la multiplication à droite par \(F_{\rm av}\) donne
\[
 \begin{pmatrix}
 F_n&F_{n-1}+F_n\\
 F_{n+1}&F_n+F_{n+1}
 \end{pmatrix}
 =
 \begin{pmatrix}
 F_n&F_{n+1}\\
 F_{n+1}&F_{n+2}
 \end{pmatrix},
\]
ce qui est le cas \(n+1\).
\end{proof}

\begin{theorem}[Identité de Cassini]
\label{p12-83-c9-s3:thm:cassini-phi}
Pour tout \(n\geq1\),
\begin{equation}
 \boxed{F_{n+1}F_{n-1}-F_n^2=(-1)^n.}
 \label{p12-83-c9-s3:eq:cassini-phi}
\end{equation}
\end{theorem}

\begin{proof}
Le déterminant de \(F_{\rm av}\) est \(-1\). En prenant les déterminants dans
\eqref{p12-83-c9-s3:eq:potencias-Fav-fibonacci},
\[
 F_{n-1}F_{n+1}-F_n^2
 =\det(F_{\rm av}^{\,n})=(-1)^n.
\]
\end{proof}

Nous définissons \(r_n=F_{n+1}/F_n\) pour \(n\geq1\). De Cassini découle
\[
 r_n-r_{n-1}
 =\frac{F_{n+1}F_{n-1}-F_n^2}{F_nF_{n-1}}
 =\frac{(-1)^n}{F_nF_{n-1}}.
\]
Par conséquent, les approximations alternent autour de la demi-droite de Perron.

\begin{proposition}[Encadrements rationnels de Perron]
\label{p12-83-c9-s3:prop:horquillas-perron-phi}
Pour \(m\geq1\),
\begin{equation}
 \frac{F_{2m+2}}{F_{2m+1}}
 <\varphi<
 \frac{F_{2m+1}}{F_{2m}},
 \label{p12-83-c9-s3:eq:horquilla-fibonacci-phi}
\end{equation}
et la largeur est
\begin{equation}
 \boxed{
 \frac{F_{2m+1}}{F_{2m}}
 -\frac{F_{2m+2}}{F_{2m+1}}
 =\frac{1}{F_{2m}F_{2m+1}}.}
 \label{p12-83-c9-s3:eq:anchura-fibonacci-phi}
\end{equation}
\end{proposition}

\begin{proof}
L’alternance précédente place les quotients pairs en dessous et les impairs
au-dessus de \(\varphi\). La formule de la largeur est
\eqref{p12-83-c9-s3:eq:cassini-phi} avec \(n=2m+1\), après réduction des deux fractions au
même dénominateur. Comme \(F_n\to\infty\), la largeur tend vers zéro.
\end{proof}

\subsection{Incidence de mémoire d’ordre \(1\) et mesure de Parry de l’enveloppe symbolique}

Soit \(X_{\rm av}\) le plus petit sous-décalage de type fini sur
\(\{\mathscr H_{\rm ar},\mathscr H_{\rm vol}\}\) qui admet exactement
les trois paires observées dans le calendrier modal. Sa matrice d’adjacence est
\(F_{\rm av}\).

\begin{theorem}[Transition et mesure de Parry]
\label{p12-83-c9-s3:thm:parry-phi}
Soit \(r=(1,\varphi)^{\mathsf T}\). La matrice stochastique de Parry est
\begin{equation}
 \boxed{
 P_{ij}=
 \frac{(F_{\rm av})_{ij}r_j}{\varphi r_i}
 =
 \begin{pmatrix}
 0&1\\
 \varphi^{-2}&\varphi^{-1}
 \end{pmatrix}.}
 \label{p12-83-c9-s3:eq:parry-matriz-phi}
\end{equation}
La mesure stationnaire des sommets est proportionnelle à
\((1,\varphi^2)\), donc
\begin{equation}
 \boxed{
 \frac{\mu_{\rm ar}}{\mu_{\rm vol}}=\varphi^{-2}.}
 \label{p12-83-c9-s3:eq:parry-razon-phi}
\end{equation}
\end{theorem}

\begin{proof}
La somme de la ligne \(i\) de \(P\) est
\[
 \frac{(F_{\rm av}r)_i}{\varphi r_i}=1.
\]
Comme \(F_{\rm av}\) est symétrique, les vecteurs propres gauche et droit
coïncident. Les poids stationnaires sont proportionnels à leurs produits
coordonnée par coordonnée, c’est-à-dire à \((1,\varphi^2)\). Le rapport entre les
deux poids est \(\varphi^{-2}\).
\end{proof}

\begin{remark}[Trois mesures à ne pas confondre]
La fréquence chronologique de l’orbite primitive est
\[
 \nu_{\rm cyc}(\mathscr H_{\rm ar})=\frac13,
 \qquad
 \nu_{\rm cyc}(\mathscr H_{\rm vol})=\frac23.
\]
La mesure uniforme \(\tau_{468}\) vit sur le catalogue des émissions. La
mesure de Parry vit sur les histoires admissibles de longueur arbitraire. La
première compte les instants, la deuxième compte les émissions et la troisième pondère
les cylindres dynamiques. Aucune n’est une approximation des autres.
\end{remark}

\subsection{Caractère multiplicatif et conservation cylindrique}

Pour une route admissible \(\gamma:i\to j\) de longueur \(n\), nous définissons
\begin{equation}
 \chi_P(\gamma)=\varphi^n\frac{r_i}{r_j}.
 \label{p12-83-c9-s3:eq:caracter-parry-phi}
\end{equation}

\begin{proposition}[Multiplicativité par concaténation]
\label{p12-83-c9-s3:prop:multiplicatividad-parry-phi}
Si \(\gamma_1:i\to j\) et \(\gamma_2:j\to k\), alors
\[
 \chi_P(\gamma_2\circ\gamma_1)
 =\chi_P(\gamma_2)\chi_P(\gamma_1).
\]
Sur une route fermée de longueur \(n\),
\[
 \chi_P(\gamma)=\varphi^n.
\]
\end{proposition}

\begin{proof}
La longueur s’additionne et la coordonnée intermédiaire se simplifie par télescopage :
\[
 \varphi^{n_1+n_2}\frac{r_i}{r_k}
 =\left(\varphi^{n_1}\frac{r_i}{r_j}\right)
  \left(\varphi^{n_2}\frac{r_j}{r_k}\right).
\]
Si \(i=k\), le quotient terminal est égal à un.
\end{proof}

Nous normalisons le vecteur gauche comme
\[
 \ell=\frac{r}{1+\varphi^2},
 \qquad
 \ell^{\mathsf T}r=1.
\]
La mesure d’un cylindre est
\[
 \mu[x_0\cdots x_n]
 =\ell_{x_0}r_{x_n}\varphi^{-n}.
\]
Par conséquent,
\begin{equation}
 V_n(x)
 :=\frac{\mu[x_0]}{\mu[x_0\cdots x_n]}
 =\varphi^n\frac{r_{x_0}}{r_{x_n}}
 =\chi_P(x_0\cdots x_n).
 \label{p12-83-c9-s3:eq:volumen-cilindrico-phi}
\end{equation}
Pour une dimension typée \(D>0\), l’échelle
\[
 a_n^{(D)}(x)=V_n(x)^{1/D}
\]
satisfait la loi exacte
\begin{equation}
 \boxed{
 \bigl(a_n^{(D)}(x)\bigr)^D
 \mu[x_0\cdots x_n]=\mu[x_0].}
 \label{p12-83-c9-s3:eq:conservacion-cilindrica-phi}
\end{equation}

En dimension trois et sur des retours fermés,
\[
 \frac{a_9}{a_0}=\varphi^3,
 \qquad
 \frac{a_{27}}{a_0}=\varphi^9,
 \qquad
 \frac{a_{108}}{a_0}=\varphi^{36}.
\]
Si \(A_k=a_{9k}\), alors
\begin{equation}
 A_{k+1}-2A_k+A_{k-1}
 =A_{k-1}(\varphi^3-1)^2>0.
 \label{p12-83-c9-s3:eq:convexidad-retornos-phi}
\end{equation}
C’est une convexité en profondeur de retour ; elle n’est pas interprétée comme une
grandeur cinématique sans une carte supplémentaire de durée.

\subsection{Mémoire quadratique et branche contractée}

Les valeurs propres de \(F_{\rm av}\) sont
\[
 \varphi,
 \qquad
 -\varphi^{-1}.
\]
La branche contractée change d’orientation à chaque itération. Pour conserver
ce signe avant d’élever au carré, nous passons à
\(\operatorname{Sym}^2(\mathbb R^2)\). Dans la base
\((x^2,2xy,y^2)\),
\begin{equation}
 \operatorname{Sym}^2(F_{\rm av})=
 \begin{pmatrix}
 0&0&1\\
 0&1&2\\
 1&1&1
 \end{pmatrix}.
 \label{p12-83-c9-s3:eq:sym2-Fav-phi}
\end{equation}

\begin{theorem}[Spectre de la mémoire de second ordre]
\label{p12-83-c9-s3:thm:espectro-sym2-phi}
\[
 \det\bigl(tI-\operatorname{Sym}^2(F_{\rm av})\bigr)
 =(t+1)(t^2-3t+1),
\]
et
\begin{equation}
 \boxed{
 \operatorname{Spec}\bigl(\operatorname{Sym}^2(F_{\rm av})\bigr)
 =\{\varphi^2,-1,\varphi^{-2}\}.}
 \label{p12-83-c9-s3:eq:espectro-sym2-phi}
\end{equation}
L’unique mode stable positif est \(\varphi^{-2}\).
\end{theorem}

\begin{proof}
Les valeurs propres d’une puissance symétrique de degré deux sont les produits
\(\lambda_1^2\), \(\lambda_1\lambda_2\) et \(\lambda_2^2\). Avec
\(\lambda_1=\varphi\) et \(\lambda_2=-\varphi^{-1}\), on obtient
\(\varphi^2\), \(-1\) et \(\varphi^{-2}\). La première dépasse un, la
deuxième a pour module un et la troisième appartient à \((0,1)\).
\end{proof}

Sur une route fermée avec \(V_n=\varphi^n\), le mode stable satisfait
\begin{equation}
 \boxed{M_n=\varphi^{-2n}=V_n^{-2}.}
 \label{p12-83-c9-s3:eq:modo-estable-phi}
\end{equation}

\subsection{Réalisation hyperbolique normalisée de l’auto-échelle : \(\exp(2\log\varphi\,J_h)=F_{\mathrm{av}}^2\)}

Soit
\[
 A=F_{\rm av}^{\,2}=
 \begin{pmatrix}1&1\\1&2\end{pmatrix},
 \qquad
 J_h=\frac{2A-3I}{\sqrt5}.
\]
Alors
\[
 J_h^2=I,
 \qquad
 A=\frac32I+\frac{\sqrt5}{2}J_h.
\]
Comme
\[
 \cosh(2\log\varphi)=\frac32,
 \qquad
 \sinh(2\log\varphi)=\frac{\sqrt5}{2},
\]
on obtient
\begin{equation}
 \boxed{
 \exp(2\log\varphi\,J_h)=F_{\rm av}^{\,2}.}
 \label{p12-83-c9-s3:eq:realizacion-hiperbolica-phi}
\end{equation}
La constante d’auto-échelle est donc simultanément la valeur propre de
Perron de l’incidence et le paramètre exponentiel de sa réalisation
hyperbolique normalisée.

\subsection{Caractérisation, encadrements et prolongement}

La construction finie détermine la demi-droite positive de \(F_{\rm av}\). La
caractérisation analytique ultérieure est
\[
 x^2-x-1=0,
 \qquad x>0.
\]
Les encadrements \eqref{p12-83-c9-s3:eq:horquilla-fibonacci-phi} ont des extrémités rationnelles
et une largeur explicite \eqref{p12-83-c9-s3:eq:anchura-fibonacci-phi}. Pour chaque niveau
nonadique \(K\), nous choisissons le plus petit \(m=m_{\varphi}(K)\) dont l’encadrement est
contenu dans un unique sous-cylindre parmi les \(729\) éléments de la fibre
ambiante. La minimalité fait de la localisation une fonction, non un
choix tacite.

Les cylindres sélectionnés sont compatibles avec les troncatures et leurs
diamètres tendent vers zéro. Leur intersection contient un unique réel ; le
\cref{p12-83-c9-s3:thm:rayo-perron-phi} l’identifie à \(\varphi\). Les quotients de
Fibonacci sont donc des certificats rationnels de la localisation, non des données
décimales employées pour décider du germe.

\subsection{Certificat interne d'unicité, de compatibilité et de convergence de l'auto-échelle}

\begin{criterion}[Réfutation de la biographie d'auto-échelle]
La construction est réfutée si l'un des faits
suivants se produit :
\begin{enumerate}
 \item le prédicat fini sélectionne plus d’une orbite ou plus d’une émission ;
 \item la récurrence \(L_0,L_0,L_1,L_1\) ne produit pas
       \eqref{p12-83-c9-s3:eq:w30-phi} ;
 \item le calendrier ne produit pas exactement les trois paires utilisées dans
       \eqref{p12-83-c9-s3:eq:N3-phi} ;
 \item un multiple de dénombrement \(N_{9},N_{27},N_{108}\) est confondu avec une
       puissance de \(N_3\) ;
 \item \(F_{\rm av}\) possède une autre demi-droite propre positive ;
 \item \eqref{p12-83-c9-s3:eq:potencias-Fav-fibonacci} ou l’identité de Cassini échoue ;
 \item les encadrements rationnels ne sont pas emboîtés ou ne contractent pas leur largeur ;
 \item la fréquence chronologique, la mesure uniforme du
       catalogue ou la mesure de Parry sont identifiées ;
 \item une table décimale intervient avant la construction de la demi-droite de Perron ;
 \item un cylindre sélectionné ne se tronque pas au cylindre précédent.
\end{enumerate}
\end{criterion}


% Unidad U015. Acoplamiento de horquillas racionales, refinamiento TPK y
% certificación a profundidad arbitraria.

\section{Certification diagonale, limite projective et génération à une précision arbitraire}
\label{p12-83-c9-s4:chap:precision-arbitraria-constantes}

\subsection{Séparation typée des \(4\) indices du prolongement nonadique}

Cette unité a le statut de certification ultérieure. Les préfixes et leurs
extensions compatibles procèdent de la relation interne
\(\operatorname{Ext}\), de la survie coinductive et des lois de
transition APP--TRIT--TPK. Les encadrements rationnels qui suivent ne choisissent pas
parmi les \(729\) éléments d’une fibre : ils localisent dans la carte archimédienne
l’extension déjà engendrée et certifient sa précision à une profondeur arbitraire.

Pour chaque caractère \(\chi\in\{\pi,e,\varphi\}\) interviennent quatre
paramètres de natures différentes :
\begin{enumerate}
 \item \(n\), ordre de l’encadrement analytique rationnel ;
 \item \(K\), niveau de raffinement du Topological Phase Kernel ;
 \item \(L_K=6K\), longueur ternaire du préfixe publié ;
 \item \(M\), nombre de positions décimales demandé.
\end{enumerate}
On ne postule ni \(n=K\), ni \(n=L_K\), ni \(K=M\), ni une proportionnalité fixe
entre eux. Le couplage est construit par des inclusions d’intervalles
décidées par arithmétique entière.

Les parties entières sont
\[
 m_\pi=3,
 \qquad
 m_e=2,
 \qquad
 m_\varphi=1,
\]
et les caractères fractionnaires sont
\[
 \rho_\pi=\pi-3,
 \qquad
 \rho_e=e-2,
 \qquad
 \rho_\varphi=\varphi-1.
\]
Les trois familles d’encadrements semi-ouverts s’écrivent
\[
 J_{\chi,n}=[\ell_{\chi,n},h_{\chi,n}).
\]
Leurs extrémités sont rationnelles, contiennent \(\rho_\chi\) et satisfont
\[
 \ell_{\chi,n}\nearrow\rho_\chi,
 \qquad
 h_{\chi,n}\searrow\rho_\chi,
 \qquad
 h_{\chi,n}-\ell_{\chi,n}\longrightarrow0.
\]
Pour \(\pi\), les extrémités proviennent des sommes alternées de l’identité
de Machin démontrée dans le \cref{chap:biografia-pi-autonoma} ; pour \(e\),
des sommes factorielles et de leur reste rationnel du
\cref{p12-83-c9-s2:chap:biografia-completa-e} ; pour \(\varphi\), des quotients de
Fibonacci et de Cassini du \cref{p12-83-c9-s3:chap:biografia-completa-phi}.

\subsection{Cylindres ternaires et fibre ambiante immédiate}

Un préfixe ternaire de longueur \(L_K=6K\) est représenté par un entier
\(N_K\) avec
\[
 0\leq N_K<3^{L_K}.
\]
Son cylindre semi-ouvert est
\begin{equation}
 C(N_K,L_K)=
 \left[
 \frac{N_K}{3^{L_K}},
 \frac{N_K+1}{3^{L_K}}
 \right).
 \label{p12-83-c9-s4:eq:cilindro-nivel-K}
\end{equation}
L’émission suivante ajoute six trits. La fibre ambiante immédiate contient
\(3^6=729\) éléments, indexés par \(u\in\{0,\ldots,728\}\) :
\begin{equation}
 C_{K,u}=
 \left[
 \frac{729N_K+u}{3^{L_K+6}},
 \frac{729N_K+u+1}{3^{L_K+6}}
 \right).
 \label{p12-83-c9-s4:eq:729-subcilindros}
\end{equation}
Les \(C_{K,u}\) sont disjoints et
\[
 C(N_K,L_K)=\bigsqcup_{u=0}^{728}C_{K,u}.
\]
Le nombre \(729\) compte les coordonnées de cette fibre ambiante ; il ne compte
ni les germes, ni les émissions distinctes, ni les orbites, ni les états survivants.

\subsubsection{Comparaison diophantienne des extrémités rationnelles de l'encadrement}

Soit \(J=[\ell,h)\subset C(N_K,L_K)\), avec
\(\ell,h\in\mathbb Q\). Nous écrivons \(S_K=3^{L_K+6}\) et définissons
\begin{align}
 j_{\min}(J,K)
 &=\max\!\left(729N_K,\left\lfloor S_K\ell\right\rfloor\right),
 \label{p12-83-c9-s4:eq:jmin-diagonal}\\
 j_{\max}(J,K)
 &=\min\!\left(729N_K+728,\left\lceil S_Kh\right\rceil-1\right).
 \label{p12-83-c9-s4:eq:jmax-diagonal}
\end{align}

\begin{lemma}[Critère de localisation unitaire]
\label{p12-83-c9-s4:lem:criterio-localizacion-unitaria}
L’encadrement \(J\) est contenu dans un unique élément de la fibre
\eqref{p12-83-c9-s4:eq:729-subcilindros} si et seulement si
\begin{equation}
 \boxed{j_{\min}(J,K)=j_{\max}(J,K).}
 \label{p12-83-c9-s4:eq:jmin-igual-jmax}
\end{equation}
Dans ce cas, si \(j\) est la valeur commune,
\[
 u=j-729N_K,
 \qquad
 N_{K+1}=j.
\]
\end{lemma}

\begin{proof}
Les entiers \(j\) dont les intervalles élémentaires de longueur \(S_K^{-1}\)
rencontrent \(J\) forment l’intervalle entier
\[
 \left[
 \lfloor S_K\ell\rfloor,
 \lceil S_Kh\rceil-1
 \right]\cap
 [729N_K,729N_K+728].
\]
Celui-ci contient exactement un élément si et seulement si ses extrémités coïncident.
L’égalité identifie le numérateur du sous-cylindre et, en tronquant ses six
derniers trits, on retrouve \(N_K\).
\end{proof}

\subsection{Certification diagonale minimale}

La génération ne fixe pas à l’avance un même ordre analytique pour tous les
niveaux. Une fois l’extension compatible produite, l’ordre de
l’encadrement n’est augmenté que jusqu’à sa certification unitaire.

\begin{definition}[Récurrence diagonale de certification]
\label{p12-83-c9-s4:def:recurrencia-diagonal-constantes}
Une fois fixée par \(\operatorname{Ext}\) une famille de cylindres compatibles
\(C(N_K,L_K)\) dont la publication contient \(\rho_\chi\), on choisit un ordre
initial \(n_\chi(K_0-1)\geq1\) et nous définissons récursivement
\begin{equation}
 n_\chi(K)=
 \min\left\{
 n\geq n_\chi(K-1):
 j_{\min}(J_{\chi,n},K)=j_{\max}(J_{\chi,n},K)
 \right\},
 \label{p12-83-c9-s4:eq:nchi-K-minimo}
\end{equation}
et nous vérifions
\begin{equation}
 N_{\chi,K+1}
 =j_{\min}(J_{\chi,n_\chi(K)},K).
 \label{p12-83-c9-s4:eq:Nchi-K-mas-uno}
\end{equation}
L’ensemble
\[
 \Delta_\chi=
 \{(n_\chi(K),K):K\geq K_0\}
\]
est la diagonale analytique--nonadique du caractère \(\chi\).
\end{definition}

\begin{theorem}[Existence et unicité de la diagonale de certification]
\label{p12-83-c9-s4:thm:existencia-unicidad-diagonal}
Pour \(\chi\in\{\pi,e,\varphi\}\), la récurrence
\eqref{p12-83-c9-s4:eq:nchi-K-minimo} est définie à tout niveau fini. L’indice
\(n_\chi(K)\) et la coordonnée \(N_{\chi,K+1}\) sont uniques.
\end{theorem}

\begin{proof}
Supposons engendrée intérieurement l’extension cylindrique de niveau \(K\) qui contient
\(\rho_\chi\). Comme \(\rho_\chi\) est irrationnel, il n’appartient à aucune
frontière ternaire \(3^{-L}\mathbb Z\). Il est donc à l’intérieur d’un
unique sous-cylindre \(C_{K,u_*}\), à une distance positive de ses deux extrémités.
Les encadrements \(J_{\chi,n}\) contiennent \(\rho_\chi\) et leur diamètre
tend vers zéro ; il existe donc un \(n\) pour lequel l’encadrement est
contenu dans \(C_{K,u_*}\). L’ensemble sur lequel est pris le minimum dans
\eqref{p12-83-c9-s4:eq:nchi-K-minimo} est non vide et, par le bon ordre de
\(\mathbb N\), possède un plus petit élément. Le
\cref{p12-83-c9-s4:lem:criterio-localizacion-unitaria} prouve que l’encadrement la
localise de manière unitaire et que le numérateur localisé coïncide avec celui déjà
engendré. La récurrence sur \(K\) complète l’argument.
\end{proof}

\begin{corollary}[Absence de paramètres décimaux]
La diagonale est calculée avec des sommes, des produits, des puissances, des divisions euclidiennes,
des parties entières inférieures et supérieures d’entiers. Elle ne reçoit aucun chiffre décimal de \(\pi\), de \(e\) ou de
\(\varphi\) comme donnée de certification.
\end{corollary}

\subsection{Compatibilité avec la connexion nonadique conjointe}

Soit \(\widehat x_{\chi,K}\) l’état enrichi du canal \(\chi\) au
niveau \(K\). Ses coordonnées incluent le préfixe, le cylindre, la phase,
le quotient, la retenue, la mémoire verticale non bornée, le terme terminal
et le registre de provenance. La phase résiduelle est
\[
 g_0(K)=[K]_9\in\mathbb Z/9\mathbb Z,
\]
tandis que l’étiquette publique est
\[
 g(K)=1+((K-1)\bmod9)\in\{1,\ldots,9\}.
\]
La mémoire verticale \(q_{\rm mem}(K)\in\mathbb Z_{\geq0}\) n’est pas réduite
modulo \(12\) ; seule l’est sa projection dodécaphasique
\[
 \beta(K)=q_{\rm mem}(K)\bmod12
\]
le fait.

Pour chaque \(K\), soit
\[
 \Gamma_{9,K}:\widehat X_K\longrightarrow\widehat X_{K+9}
\]
l’holonomie d’un tour de phase. Alors
\[
 g_0(K+9)=g_0(K),
 \qquad
 q_{\rm mem}(K+9)=q_{\rm mem}(K)+1,
 \qquad
 \beta(K+9)=\beta(K)+1\pmod{12}.
\]

\begin{theorem}[Retour de phase avec conservation généalogique]
\label{p12-83-c9-s4:thm:retorno-fase-conservacion-genealogica}
L’égalité de phase ne réinitialise pas l’état :
\[
 g_0(K+9)=g_0(K),
 \qquad
 \widehat x_{\chi,K+9}\ne\widehat x_{\chi,K}.
\]
Les profondeurs \(27\), \(54\) et \(108\) sont les compositions typées
\begin{align*}
 \Gamma_{27,K}
 &=\Gamma_{9,K+18}\circ\Gamma_{9,K+9}\circ\Gamma_{9,K},\\
 \Gamma_{54,K}
 &=\Gamma_{9,K+45}\circ\cdots\circ\Gamma_{9,K},\\
 \Gamma_{108,K}
 &=\Gamma_{9,K+99}\circ\cdots\circ\Gamma_{9,K}.
\end{align*}
En aucun cas il ne s’agit de réinitialisations du préfixe, de la retenue, du terminal ou de la
mémoire.
\end{theorem}

\begin{proof}
La première coordonnée est calculée dans \(\mathbb Z/9\mathbb Z\), tandis que
\(q_{\rm mem}\) appartient à un monoïde non borné. Chaque application de
\(\Gamma_{9,K}\) ajoute une unité de mémoire, six blocs par niveau
intermédiaire et la composition correspondante des cocycles. C’est pourquoi la phase
peut coïncider sans que l’état enrichi coïncide. Les formules de
profondeur montrent les indices de chaque application et évitent de les composer comme si toutes
étaient des endomorphismes du même espace.
\end{proof}

\subsubsection{Compression, expression et terme terminal}

La réalisation opératorielle d’un tour satisfait
\begin{equation}
 U_\Gamma J=JT+\eta,
 \qquad
 J^*\eta=0,
 \qquad
 I-T^*T=\eta^*\eta.
 \label{p12-83-c9-s4:eq:compresion-expresion-precision}
\end{equation}
La contraction visible stricte est réservée à
\[
 M_9=\frac59V_9;
\]
n’est pas identifié à l’opérateur de compression \(T\). Le télescopage de
profondeur \(L\) conserve le terme terminal :
\begin{equation}
 S_0=
 \sum_{r=0}^{L-1}M_9^{*r}D_9^*D_9M_9^r
 +M_9^{*L}S_0M_9^L.
 \label{p12-83-c9-s4:eq:telescopia-terminal-precision}
\end{equation}
Omettre le dernier terme avant de prouver sa disparition effacerait
de l’information de frontière et transformerait le retour de phase en une fausse
réinitialisation.

\subsection{Système inverse des histoires compatibles}

Soit \(\mathcal H_K^\chi\) l’ensemble des histoires enrichies admissibles
du canal \(\chi\) à la profondeur \(K\), et soit
\[
 p_{N,K}:\mathcal H_N^\chi\longrightarrow\mathcal H_K^\chi
 \qquad(N\geq K)
\]
la troncature. Nous définissons
\[
 \operatorname{Surv}_K^{(N)}(\chi)
 =p_{N,K}(\mathcal H_N^\chi)
\]
et
\begin{equation}
 \operatorname{Surv}_K(\chi)
 =\bigcap_{N\geq K}\operatorname{Surv}_K^{(N)}(\chi).
 \label{p12-83-c9-s4:eq:supervivencia-todas-profundidades}
\end{equation}
L’objet coinductif est
\begin{equation}
 X_\chi=
 \varprojlim_K
 \bigl(\operatorname{Surv}_K(\chi),p_{K+1,K}\bigr).
 \label{p12-83-c9-s4:eq:limite-inverso-caracteres}
\end{equation}

\begin{theorem}[Section compatible sous prolongeabilité enrichie]
\label{p12-83-c9-s4:thm:seccion-compatible-caracteres}
Supposons en outre que, sur chaque sous-cylindre numérique engendré, la
relation \(\operatorname{Ext}\) soit non vide, univoque et compatible avec les
troncatures. Alors \(\operatorname{Ext}\) et la survie produisent une famille
\[
 x_\chi=(\widehat x_{\chi,K})_{K\geq K_0}
\]
qui satisfait
\[
 p_{K+1,K}(\widehat x_{\chi,K+1})
 =\widehat x_{\chi,K}.
\]
Par conséquent, \(x_\chi\in X_\chi\).
\end{theorem}

\begin{proof}
Le sous-cylindre de niveau \(K+1\) a pour numérateur
\(N_{\chi,K+1}=729N_{\chi,K}+u\). La troncature élimine ses six derniers
trits et retrouve \(N_{\chi,K}\). Les autres coordonnées sont obtenues par les
lois fonctorielles de retenue, de mémoire, de terminal et de phase. L’hypothèse
d’univocité de \(\operatorname{Ext}\) fixe le sous-cylindre et l’extension
enrichie ; l’unicité locale du
\cref{p12-83-c9-s4:lem:criterio-localizacion-unitaria} certifie ensuite sa
localisation archimédienne, et la compatibilité de \(\operatorname{Ext}\)
fournit l’égalité de troncature.
\end{proof}

\subsection{Cellules décimales et précision prescrite}

Pour \(M\geq1\) et \(q\in\{0,\ldots,10^M-1\}\), soit
\[
 D_{M,q}=
 \left[
 \frac{q}{10^M},
 \frac{q+1}{10^M}
 \right).
\]
Un encadrement \(J=[\ell,h)\) est contenu dans une unique cellule décimale si
et seulement si
\begin{equation}
 \boxed{
 \left\lfloor10^M\ell\right\rfloor
 =\left\lceil10^Mh\right\rceil-1=q.}
 \label{p12-83-c9-s4:eq:criterio-celda-decimal}
\end{equation}

\begin{definition}[Niveau minimal de certification décimale]
Pour \(\chi\in\{\pi,e,\varphi\}\), nous définissons \(K_\chi(M)\) comme le plus petit
niveau pour lequel il existe un ordre analytique \(n\) et un entier \(q\) tels
que
\begin{enumerate}
 \item \(J_{\chi,n}\) satisfait
       \eqref{p12-83-c9-s4:eq:criterio-celda-decimal} ;
 \item \(J_{\chi,n}\) est contenu dans le cylindre engendré par
       l’état \(\widehat x_{\chi,K}\) ;
 \item la transition vers le niveau suivant satisfait
       \eqref{p12-83-c9-s4:eq:jmin-igual-jmax}.
\end{enumerate}
La paire \((K,n)\) est minimisée lexicographiquement.
\end{definition}

\begin{theorem}[Certification à précision décimale arbitraire des sorties coinductives]
\label{p12-83-c9-s4:thm:precision-decimal-arbitraria}
Pour chaque \(\chi\in\{\pi,e,\varphi\}\) et chaque \(M\geq1\), il existe
\(K_\chi(M)<\infty\). La section engendrée détermine un unique préfixe décimal
de \(M\) positions, et l’encadrement reconnaît et certifie ce préfixe. Si
\(M'<M\), le préfixe d’ordre \(M\) se tronque à celui
d’ordre \(M'\), et les cellules correspondantes sont emboîtées.
\end{theorem}

\begin{proof}
Les trois caractères sont irrationnels : \(\varphi\) l’est par son polynôme
quadratique à discriminant non carré ; l’irrationalité de \(e\) et de
\(\pi\) est tirée de leurs théorèmes classiques de caractérisation, postérieurs à
la construction finie. Aucun n’appartient donc à une frontière décimale ou
ternaire rationnelle. Pour un \(M\) fixé, la distance du caractère à l’ensemble
fini des frontières décimales pertinentes est positive. Les encadrements
rationnels contractent leur diamètre jusqu’à être contenus dans une unique cellule décimale.
Le \cref{p12-83-c9-s4:thm:existencia-unicidad-diagonal} couple cet encadrement à un
cylindre TPK et à une extension unique. Le minimum lexicographique existe par
bon ordre. La compatibilité de la section prouve l’emboîtement lorsque
\(M\) augmente.
\end{proof}

\begin{corollary}[Unicité réelle]
Les cellules emboîtées ont des diamètres \(10^{-M}\to0\). Leur intersection est
un singleton, identifié par les caractérisations précédentes à
\(\pi\), \(e\) ou \(\varphi\), respectivement.
\end{corollary}

\subsection{Équivalence de \(2\) réalisations finies de la loi uniforme}

Les témoins de \(1,000\) et de \(10,000\) positions ne définissent pas les
caractères et ne limitent pas le prolongement. Ce sont des exécutions de la même récurrence
uniforme à deux profondeurs différentes.

\begin{table}[htbp]
\centering
\caption{Dimensions exactes des exécutions certifiées, par canal.}
\label{p12-83-c9-s4:tab:dimensiones-ejecuciones-1000-10000}
\begin{tabular}{lrr}
\toprule
Objet compté & \(1,000\) positions & \(10,000\) positions\\
\midrule
Partitions exhaustives de la fibre ambiante & 396 & 3501\\
Trits transportados & 2376 & 21006\\
Paires de même phase \((K,K+9)\) & 387 & 3492\\
Tours nonadiques complets & 43 & 388\\
Éléments examinés par partition & 729 & 729\\
Extensions compatibles conservées par partition & 1 & 1\\
\bottomrule
\end{tabular}
\end{table}

Para \(10,000\) posiciones,
\[
 B=3501=389\cdot9,
 \qquad
 6B=21006,
 \qquad
 3^{6B}>10^{10000}.
\]
Les \(396\) premiers blocs de l’exécution longue coïncident avec
l’exécution de \(1,000\) positions. Dans chacune des \(3492\) paires de
même phase, la coordonnée de phase est préservée et le préfixe, la
profondeur, le cylindre et la mémoire sont prolongés. Le certificat enregistre aussi les
répétitions éventuelles de projections locales sans les confondre avec un retour
de l’état complet.

\subsection{Séparation entre génération et vérification}

\begin{definition}[Générateur structurel]
Le générateur reçoit le catalogue fini, \(L_0,L_1\), l’état initial, les
lois de connexion et la relation interne \(\operatorname{Ext}\). À chaque niveau :
\begin{enumerate}
 \item il énumère les \(729\) éléments de la fibre ambiante ;
 \item il applique \(\operatorname{Ext}\) et conserve l’unique extension compatible ;
 \item il actualise la phase, le quotient, la retenue, la mémoire et le terminal ;
 \item il écrit le bloc ternaire, le reçu de sélection et le registre généalogique.
\end{enumerate}
\end{definition}

\begin{definition}[Vérification ultérieure indépendante]
La vérification ultérieure indépendante reçoit les sorties closes du
générateur et les trois familles d'encadrements rationnels ; elle calcule
\(j_{\min}\) et \(j_{\max}\), certifie leur égalité avec l'extension déjà
produite, reconstruit les
partitions, vérifie les carrés de troncature, les couples
\((K,K+9)\), les sommes de partition
\[
 \operatorname{killed}_{\mathrm{izq}}+1+\operatorname{killed}_{\mathrm{der}}=729,
\]
l'intégrité des sorties et, à la fin, la concordance avec les développements de référence.
\end{definition}

\begin{theorem}[Isolement des données cibles]
\label{p12-83-c9-s4:thm:aislamiento-datos-objetivo-precision}
Les développements décimaux de référence et les données métrologiques
n’interviennent pas dans la sélection d’une extension compatible. La comparaison est
effectuée après l’écriture des trois préfixes.
\end{theorem}

\begin{proof}
L'entrée mathématique de chaque décision générative consiste en le catalogue
fini, \(L_0,L_1\), l'état hérité et la relation interne
\(\operatorname{Ext}\). Les extrémités rationnelles, puissances de \(3\), planchers et
plafonds relèvent de la vérification ultérieure. Le générateur termine avant
que celle-ci n'ouvre les encadrements ou n'exécute les contrôles décimaux. Le sens
du flux de données est
\[
 \text{structure}\longrightarrow
 \text{préfixes et registres}\longrightarrow
 \text{vérification ultérieure indépendante};
\]
il n’existe pas de flèche en sens contraire.
\end{proof}

\subsection{Obstructions différenciées de la génération et de sa certification à précision arbitraire}

\begin{criterion}[Réfutation de la génération à une précision arbitraire]
Les défaillances d'orbite, d'extension, de troncature ou de retour réfutent la génération ;
les défaillances d'encadrement, d'ordre ou de contrôle décimal réfutent le certificat
analytique correspondant ; et une divergence entre les deux réfute
l'identification conjointe de la sortie publiée. Avec cette séparation des types,
l'affirmation complète est réfutée si :
\begin{enumerate}
 \item on identifie l’un quelconque des indices \(n,K,L_K,M\) sans un
       théorème de liaison ;
 \item un encadrement cesse de contenir son caractère ou ne contracte pas son diamètre ;
 \item une étape satisfait \(j_{\min}\ne j_{\max}\) après l’ordre
       déclaré minimal ;
 \item une extension ne se tronque pas au cylindre du niveau précédent ;
 \item le retour de phase réinitialise le préfixe, la retenue, le terminal ou la mémoire ;
 \item le terme terminal est omis dans
       \eqref{p12-83-c9-s4:eq:telescopia-terminal-precision} sans que son annulation soit démontrée ;
 \item les \(396\) premiers blocs de l’exécution longue diffèrent de
       l’exécution courte ;
 \item un contrôle décimal ou métrologique est lu avant la clôture de la sortie ;
 \item un chiffre imprimé diffère de la section mathématique certifiée ;
 \item la preuve dépend du fait que la précision demandée soit \(1,000\) ou
       \(10,000\), au lieu d’un \(M\) arbitraire.
\end{enumerate}
\end{criterion}
